\section{Вокруг языка C++}

\subsection{Какие стандарты языка существуют? Каков цикл выпуска новых стандартов?}

Язык C++ стандартизуется. У него есть последовательность стандартов, в которых фиксируются возможности языка и стандартной библиотеки.

\begin{definition}[Основные стандарты C++]
Обычно выделяют следующие крупные ревизии языка:
\begin{itemize}
    \item \textbf{C++98};
    \item \textbf{C++03};
    \item \textbf{C++11};
    \item \textbf{C++14};
    \item \textbf{C++17};
    \item \textbf{C++20};
    \item \textbf{C++23}.
\end{itemize}
\end{definition}

Для курса важно помнить отдельно:
\begin{itemize}
    \item \textbf{базовый стандарт курса --- C++17};
    \item \textbf{отдельные элементы C++20} тоже используются и обсуждаются.
\end{itemize}

До окончательного принятия нового стандарта существует его рабочий проект --- \textit{draft}.

\begin{remark}
В современный период новые стандарты C++ выходят примерно \textbf{раз в три года}. На защите полезно разделять:
\begin{itemize}
    \item общую историю стандартов;
    \item конкретную договорённость курса: C++17 + отдельные элементы C++20.
\end{itemize}
\end{remark}


\subsection{Что такое The Zero Overhead Principle? Привести пример появления этого принципа в коде.}

\begin{definition}[The Zero Overhead Principle]
Принцип нулевых накладных расходов означает: если программист не использует некоторую абстракцию, программа не должна за неё ``платить''; если использует, цена этой абстракции должна быть минимальной.
\end{definition}

Это одна из ключевых идей C++: язык даёт высокоуровневые средства, но старается не превращать их в обязательные лишние runtime-расходы.

\begin{example}
Простой пример --- вызов небольшой функции или шаблонной функции:
\begin{verbatim}
template <class T>
T myMax(const T& a, const T& b)
{
    return (a < b) ? b : a;
}
\end{verbatim}
С точки зрения программиста это удобная абстракция. После инстанцирования и оптимизации компилятор обычно получает код, который по стоимости близок к ручной реализации для конкретного типа.
\end{example}

\begin{remark}
На этом вопросе не нужно уходить в слишком глубокую теорию оптимизаций. Важно правильно объяснить сам принцип.
\end{remark}


\subsection{Какие основные директивы языка CMake/файлов проекта? Для чего нужна директива \texttt{add\_executable}? Каковы её параметры?}

В курсе все проекты создаются на базе CMake.

\begin{definition}[Типичные команды CMake]
В простом проекте чаще всего встречаются:
\begin{itemize}
    \item \texttt{cmake\_minimum\_required(...)} --- минимальная версия CMake;
    \item \texttt{project(...)} --- имя проекта и язык;
    \item \texttt{add\_executable(...)} --- создание исполняемого таргета;
    \item \texttt{add\_library(...)} --- создание библиотечного таргета.
\end{itemize}
\end{definition}

\begin{definition}[\texttt{add\_executable}]
Команда \texttt{add\_executable} создаёт \textbf{target}, результатом сборки которого является исполняемый файл.
\end{definition}

В простейшем случае у неё два главных смысловых параметра:
\begin{itemize}
    \item имя таргета;
    \item список исходных файлов.
\end{itemize}

\begin{example}
\begin{verbatim}
cmake_minimum_required(VERSION 3.5)
project(WorkshopX CXX)

add_executable(ex1 ex1/main.cpp)
\end{verbatim}
\end{example}

\begin{remark}
На защите лучше использовать слово \textbf{target}, потому что именно этот термин используется в лекциях.
\end{remark}


\subsection{Что такое интерфейс и что такое реализация? Что из них можно отнести к идеальному, а что --- к материальному?}

\begin{definition}[Интерфейс]
Интерфейс объявляет что-то без задания конкретного способа реализации и тем самым задаёт \textbf{контракт} между частями программы.
\end{definition}

\begin{definition}[Реализация]
Реализация претворяет этот контракт в жизнь одним из возможных способов.
\end{definition}

Интерфейс отвечает на вопрос \emph{что можно делать}, а реализация --- \emph{как это устроено}.

Одному и тому же интерфейсу может соответствовать несколько разных реализаций.

С идейной точки зрения:
\begin{itemize}
    \item \textbf{интерфейс} относится к более \textbf{идеальной} стороне;
    \item \textbf{реализация} относится к более \textbf{материальной} стороне.
\end{itemize}


\subsection{Привести конкретные примеры разделения на интерфейс и реализацию в языке C++ (модули, типы и т.д.).}

Типичные примеры в C++:

\begin{itemize}
    \item \textbf{Прототип функции} --- интерфейс функции.
    \item \textbf{Тело функции} --- реализация.
    \item \textbf{Создание нового типа} --- часть интерфейсного описания.
    \item \textbf{Создание объекта конкретного типа} --- уже реализационная часть.
    \item \textbf{Интерфейсный класс} --- класс, состоящий только из чисто виртуальных функций.
    \item \textbf{Наследование от интерфейсного класса и override виртуальных функций} --- реализация интерфейса.
\end{itemize}

\begin{example}
\begin{verbatim}
double power(double x, unsigned int p);   // интерфейс

double power(double x, unsigned int p)
{
    double res = x;
    for (unsigned int i = 1; i < p; ++i)
        res = res * x;
    return res;
}
\end{verbatim}
\end{example}


\subsection{Для чего используются заголовочные файлы в C++? Что в них обычно помещают? Что с ними обычно делают и что с ними обычно делать нельзя? Какие расширения встречаются у заголовочных файлов? Какая связь между заголовочными файлами и принципом интерфейса-реализации?}

\begin{definition}[Заголовочный файл]
Заголовочный файл --- это, как правило, \textbf{интерфейсная часть} некоторого модуля.
\end{definition}

Обычно в заголовочные файлы кладут:
\begin{itemize}
    \item прототипы функций;
    \item описания пользовательских типов;
    \item интерфейсные классы;
    \item в ряде случаев --- шаблоны целиком.
\end{itemize}

Типичные расширения:
\begin{itemize}
    \item \texttt{.h};
    \item \texttt{.hpp};
    \item \texttt{.hxx}.
\end{itemize}

Что с ними обычно делают:
\begin{itemize}
    \item подключают через \texttt{\#include} в \texttt{.cpp};
    \item через них распространяют интерфейс между модулями.
\end{itemize}

Что обычно делать не стоит:
\begin{itemize}
    \item помещать туда обычные определения объектов и функций, если заголовок будет включаться более чем в один модуль трансляции;
    \item класть туда детали реализации, не относящиеся к интерфейсу.
\end{itemize}

\begin{remark}
Заголовочный файл в обычной организации проекта соответствует именно \textbf{интерфейсу}.
\end{remark}


\subsection{Для чего используются модули трансляции в C++? Что в них обычно помещают? Что с ними обычно делают и что с ними обычно делать нельзя? Какие расширения встречаются у модулей трансляции? Какая связь между модулями трансляции и принципом интерфейса-реализации?}

\begin{definition}[Модуль трансляции]
Модуль трансляции (\textit{translation unit}) --- это исходный файл C++, который после работы препроцессора компилируется как единое целое.
\end{definition}

В рамках обычной организации проекта это, как правило, \texttt{.cpp}-файл.

Обычно в модуле трансляции помещают:
\begin{itemize}
    \item реализации функций;
    \item определения объектов;
    \item детали реализации, скрытые от внешнего кода.
\end{itemize}

Что с ними делают:
\begin{itemize}
    \item препроцессируют;
    \item компилируют;
    \item затем линкуют результат в программу.
\end{itemize}

Типичное расширение:
\begin{itemize}
    \item \texttt{.cpp}.
\end{itemize}

Важно различать:
\begin{itemize}
    \item \texttt{.cpp} --- отдельный модуль трансляции;
    \item \texttt{.h}, \texttt{.hpp}, \texttt{.hxx} --- заголовок, а не отдельный компилируемый модуль.
\end{itemize}

\begin{remark}
Если заголовок --- это обычно интерфейс, то модуль трансляции --- это обычно \textbf{реализация}.
\end{remark}



\section{Основы C++}

\subsection{Каково назначение препроцессора языка C(++)? С какими сущностями он манипулирует? Имеет ли для него принципиальное значение синтаксис и семантика языка C++?}

\begin{definition}[Препроцессор]
Препроцессор --- этап обработки текста программы \emph{до} компиляции.
\end{definition}

Его назначение:
\begin{itemize}
    \item подключать заголовочные файлы;
    \item разворачивать макросы;
    \item выполнять условное включение или исключение фрагментов текста;
    \item конфигурировать текст программы под разные режимы сборки, ОС и платформы.
\end{itemize}

Препроцессор работает прежде всего с \textbf{текстом} исходника. Он манипулирует:
\begin{itemize}
    \item директивами;
    \item токенами;
    \item текстовыми фрагментами.
\end{itemize}

Для него не имеет принципиального значения полная семантика C++: он не понимает типы, классы, перегрузку и т.д. Его работа носит текстовый характер.

\begin{remark}
Полезная формулировка: препроцессор подготавливает текст исходного кода под конкретные условия билда.
\end{remark}


\subsection{На примере директивы \texttt{assert} покажите, как препроцессор меняет код программы для конфигураций сборки debug и release.}

\texttt{assert} подключается через:
\begin{verbatim}
#include <cassert>
\end{verbatim}

\begin{example}
\begin{verbatim}
#include <cassert>

int f(int x)
{
    assert(x != 0);
    return 100 / x;
}
\end{verbatim}
\end{example}

Если макрос \texttt{NDEBUG} \textbf{не определён}, проверка остаётся в программе и при ложном условии завершает выполнение с диагностикой.

Если \texttt{NDEBUG} \textbf{определён}, макрос \texttt{assert(...)} разворачивается в пустую безопасную конструкцию вида \texttt{((void)0)} и фактически исчезает из runtime-поведения программы.

Это пример того, как один и тот же текст исходника конфигурируется по-разному для debug и release.


\subsection{Что содержит и что не содержит листинг (результат работы препроцессора)?}

\begin{definition}[Листинг]
Листинг --- это текст программы после работы препроцессора.
\end{definition}

Листинг обычно содержит:
\begin{itemize}
    \item результаты \texttt{\#include};
    \item раскрытые макросы;
    \item уже выбранные ветки условной компиляции;
    \item служебные строки, помогающие соотносить результат с исходными файлами.
\end{itemize}

Листинг обычно не содержит в исходном виде:
\begin{itemize}
    \item самих директив \texttt{\#include}, \texttt{\#define} как управляющих конструкций;
    \item комментариев как самостоятельной сущности;
    \item результата семантического анализа программы компилятором.
\end{itemize}


\subsection{Приведите пример конфигурирования текста исходного кода программы условными директивами п. для... (например, включения фрагмента кода, релевантного определенной ОС).}

\begin{example}
\begin{verbatim}
#ifdef _WIN32
    // код только для Windows
#elif defined(__linux__)
    // код только для Linux
#else
    // код для остальных систем
#endif
\end{verbatim}
\end{example}

Другой частый пример --- условное включение отладочного кода:
\begin{verbatim}
#ifdef DEBUG
    std::cerr << "debug mode\n";
#endif
\end{verbatim}


\subsection{Как включить в модуль системный заголовочный файл? ---//--- заголовочный файл из проекта?}

Системный заголовок подключают так:
\begin{verbatim}
#include <iostream>
\end{verbatim}

Проектный заголовок подключают так:
\begin{verbatim}
#include "foo.h"
\end{verbatim}

Разница в синтаксисе отражает различный способ поиска файла:
\begin{itemize}
    \item \texttt{<...>} --- для системных путей;
    \item \texttt{"..."} --- сначала локальный поиск, потом системные пути.
\end{itemize}


\subsection{В чем заключается проблема повторного косвенного включения заголовочного ф. Привести пример.}

Проблема возникает, когда один и тот же заголовочный файл включается в один и тот же модуль трансляции несколько раз через цепочку других заголовков.

\begin{example}
\begin{verbatim}
// foo.h
#ifndef FOO_H_
#define FOO_H_

double power(double x, unsigned int p);
int GLOBAL_I;

#endif

// bar.h
#ifndef BAR_H_
#define BAR_H_

#include "foo.h"
double power2(double x, double p);

#endif

// main.cpp
#include "foo.h"
#include "bar.h"
\end{verbatim}
\end{example}

Даже если повторное текстовое включение самого файла уже предотвращено, в заголовке всё равно может остаться ``мина'' --- например, \texttt{int GLOBAL\_I;}. Такой заголовок можно включить только в один модуль трансляции, иначе будут проблемы на этапе линковки.


\subsection{Что такое header guard и какую проблему он призван предотвратить? Привести конкретный пример.}

\begin{definition}[Header guard]
Header guard --- это защита заголовочного файла от повторного включения.
\end{definition}

Классическая реализация:
\begin{verbatim}
#ifndef PROJECT_MATH_VECTOR_H
#define PROJECT_MATH_VECTOR_H

// содержимое заголовка

#endif
\end{verbatim}

Header guard предотвращает повторное текстовое включение содержимого одного и того же заголовка в один модуль трансляции.


\subsection{Можно ли использовать заголовочный ф. без h.g. и если да, то какие ограничения это накладывает?}

Теоретически можно, но только при очень жёстком ограничении: такой заголовок должен попасть в каждый модуль трансляции не более одного раза.

На практике это неудобно и небезопасно. Поэтому нормальная практика --- использовать либо классический header guard, либо \texttt{\#pragma once}.


\subsection{Что обычно кладут в заголовочный файл --- интерфейсную часть некоторого модуля foo? А что обычно не кладут --- и почему?}

Обычно в заголовочный файл кладут:
\begin{itemize}
    \item прототипы функций;
    \item описания типов;
    \item объявления объектов через \texttt{extern}, если это действительно нужно;
    \item шаблоны.
\end{itemize}

Обычно не кладут:
\begin{itemize}
    \item определения обычных объектов;
    \item определения обычных нешаблонных функций;
    \item детали реализации, не относящиеся к интерфейсу.
\end{itemize}

Причина --- при включении такого заголовка в несколько модулей трансляции можно получить ошибки компиляции или линковки.


\subsection{Чем в C++ описание чего-л. принципиально отличается от определения чего-л.? Как можно наблюдать эффект описания чего-л. и эффект определения чего-л.?}

\begin{definition}[Описание]
Описание (\textit{declaration}) создаёт для компилятора знание о том, что в программе есть некоторая сущность и как с ней можно работать.
\end{definition}

\begin{definition}[Определение]
Определение (\textit{definition}) создаёт реальный объект в памяти, либо даёт тело функции, то есть реальную сущность в программе.
\end{definition}

Простейшая формулировка:
\begin{itemize}
    \item \textbf{описать} --- значит оповестить компилятор;
    \item \textbf{определить} --- значит реально создать сущность.
\end{itemize}

\begin{example}
\begin{verbatim}
extern int x;   // только описание
int x = 5;      // определение и одновременно описание
\end{verbatim}
\end{example}

Для функций:
\begin{verbatim}
int sum(int, int);        // описание
int sum(int a, int b)     // определение
{
    return a + b;
}
\end{verbatim}

Эффект описания наблюдается в том, что компилятор ``знает имя'' сущности. Эффект определения наблюдается в том, что сущность начинает реально существовать в программе.



\section{Базовые конструкции C++}

\subsection{Алфавит языка C++. Простая программа C++, структура, основные действующие лица}

Алфавит C++ включает:
\begin{itemize}
    \item латинские буквы;
    \item цифры;
    \item знак подчёркивания;
    \item специальные символы операторов и разделителей;
    \item символы препроцессора;
    \item символы пробельных разделителей.
\end{itemize}

Минимальная программа:
\begin{verbatim}
int main()
{
    return 0;
}
\end{verbatim}

Основные ``действующие лица'' при построении программы:
\begin{itemize}
    \item препроцессор;
    \item компилятор;
    \item линкер;
    \item runtime-среда.
\end{itemize}


\subsection{Ключевые слова. Основные ключевые слова}

Ключевое слово --- зарезервированное слово языка, которое нельзя использовать как обычный идентификатор.

Примеры:
\begin{itemize}
    \item типы и объявления: \texttt{int}, \texttt{double}, \texttt{char}, \texttt{void}, \texttt{class}, \texttt{struct}, \texttt{enum};
    \item управление потоком: \texttt{if}, \texttt{else}, \texttt{switch}, \texttt{for}, \texttt{while}, \texttt{do}, \texttt{break}, \texttt{continue}, \texttt{return};
    \item спецификаторы: \texttt{const}, \texttt{static}, \texttt{extern}, \texttt{inline}, \texttt{virtual}, \texttt{namespace}, \texttt{auto}.
\end{itemize}


\subsection{Идентификаторы, правило составления идентификаторов. Влияние кодстайла на п.с.и.}

Идентификатор:
\begin{itemize}
    \item состоит из букв, цифр и символа \texttt{\_};
    \item не может начинаться с цифры;
    \item чувствителен к регистру;
    \item не должен совпадать с ключевым словом.
\end{itemize}

Кодстайл накладывает дополнительные соглашения:
\begin{itemize}
    \item типы часто именуют в \texttt{PascalCase};
    \item функции и переменные --- в \texttt{camelCase};
    \item макросы --- в \texttt{UPPER\_CASE};
    \item глобальные константы тоже часто выделяют отдельным стилем.
\end{itemize}


\subsection{Что такое инструкция? Какой признак в программе явно указывается на появление инструкции? Когда это не так?}

\begin{definition}[Инструкция]
Инструкция (\textit{statement}) --- законченная синтаксическая конструкция языка.
\end{definition}

Часто признаком инструкции служит:
\begin{itemize}
    \item точка с запятой \texttt{;};
    \item либо фигурные скобки для составной инструкции.
\end{itemize}

Но это не всегда так:
\begin{itemize}
    \item \texttt{if}, \texttt{switch}, циклы не обязательно заканчиваются \texttt{;} как целое;
    \item блок \texttt{\{...\}} сам является составной инструкцией.
\end{itemize}


\subsection{Что такое инструкция описания объекта? Приведите пример.}

Это инструкция, которая сообщает компилятору о существовании объекта и его типе, но не обязательно создаёт его.

\begin{example}
\begin{verbatim}
extern int counter;
\end{verbatim}
\end{example}


\subsection{Что такое инструкция описания функции? Приведите пример.}

Это инструкция, которая задаёт интерфейс функции без тела.

\begin{example}
\begin{verbatim}
double power(double x, unsigned int p);
\end{verbatim}
\end{example}

Такой заголовок функции вместе с \texttt{;} называется прототипом.


\subsection{Что такое инструкция определения и описания объекта? Приведите пример.}

Это инструкция, которая одновременно и описывает объект для компилятора, и реально создаёт его.

\begin{example}
\begin{verbatim}
int a = 42;
\end{verbatim}
\end{example}


\subsection{Возможно ли существование инструкции только описания без определения и наоборот? Приведите примеры.}

Только описание без определения --- возможно:
\begin{verbatim}
extern int x;
int sum(int, int);
\end{verbatim}

Определение без отдельной инструкции описания тоже возможно, потому что определение обычно одновременно является и описанием:
\begin{verbatim}
int x = 5;
int sum(int a, int b) { return a + b; }
\end{verbatim}


\subsection{Каким образом в C++ производится инициализация объектов и когда? Приведите различные подходы для инициализации, что между ними общего и какие есть различия.}

Инициализация происходит \textbf{в момент создания объекта}.

Основные формы:
\begin{itemize}
    \item \texttt{T x = value;} --- copy-initialization;
    \item \texttt{T x(value);} --- direct-initialization;
    \item \texttt{T x\{value\};} --- list-initialization;
    \item \texttt{T x = \{value\};} --- copy-list-initialization.
\end{itemize}

Общее: все они задают начальное значение создаваемому объекту.

Различия:
\begin{itemize}
    \item \texttt{\{\}} лучше обнаруживает narrowing;
    \item \texttt{=} в объявлении выглядит как присваивание, но здесь это именно инициализация;
    \item \texttt{()} и \texttt{\{\}} различаются правилами выбора конструкторов и поведением при сужении.
\end{itemize}


\subsection{Что такое блок инструкций? Можно ли использовать блок без каких-либо заголовочных инструкций перед ним?}

Блок инструкций --- это последовательность инструкций в фигурных скобках:
\begin{verbatim}
{
    int x = 1;
    int y = x + 2;
}
\end{verbatim}

Да, блок можно использовать и без заголовка \texttt{if}/\texttt{for}/функции. Это создаёт отдельную локальную область видимости.


\subsection{Что такое область видимости имени (scope)? Какой scope задаётся блоком инструкций?}

Scope --- область программы, внутри которой имя видно и доступно без дополнительной квалификации.

Блок инструкций задаёт \textbf{локальную область видимости}.


\subsection{Что такое функция с точки зрения языка C++?}

Функция --- именованный блок инструкций с заданным интерфейсом:
\begin{itemize}
    \item имя;
    \item возвращаемый тип;
    \item список параметров;
    \item тело.
\end{itemize}

Функция задаёт контракт вызова и реализует некоторое действие или вычисление.


\subsection{Объекты и типы. Характеристики объектов и характеристики типов.}

Объект --- конкретная область памяти с данными.

Тип --- описание множества значений и множества операций над ними.

Характеристики объекта:
\begin{itemize}
    \item тип;
    \item значение;
    \item адрес;
    \item размер;
    \item время жизни;
    \item имя, если оно есть.
\end{itemize}

Характеристики типа:
\begin{itemize}
    \item множество значений;
    \item набор операций;
    \item размер;
    \item способ кодирования;
    \item структура представления.
\end{itemize}


\subsection{Представление объектов в оперативной памяти (footprint/dump), влияние на представления размера объекта и способа его кодирования. Оператор sizeof().}

Объект реально располагается в памяти и имеет некоторый \textit{footprint} --- занимаемый объём и структуру байтов.

На представление влияют:
\begin{itemize}
    \item размер типа;
    \item способ кодирования данных;
    \item выравнивание;
    \item внутренняя структура типа.
\end{itemize}

Оператор \texttt{sizeof} возвращает размер типа или объекта в байтах:
\begin{verbatim}
sizeof(int)
sizeof x
\end{verbatim}


\subsection{Операционная семантика объекта определенного типа (операции). Абстрактный тип данных.}

Операционная семантика типа --- это набор допустимых операций над объектами данного типа и смысл этих операций.

Например, для \texttt{int} это арифметика, сравнение, присваивание, инкремент и т.д.

Абстрактный тип данных --- это тип, для которого важнее интерфейс операций, чем конкретная реализация хранения.


\subsection{Классификация типов в C++ (трехтаксонная) и по cppreference (двухтаксонная).}

Трёхтаксонная классификация:
\begin{itemize}
    \item фундаментальные;
    \item производные;
    \item пользовательские.
\end{itemize}

Примеры:
\begin{itemize}
    \item фундаментальные: \texttt{int}, \texttt{double}, \texttt{char}, \texttt{bool}, \texttt{void};
    \item производные: указатели, ссылки, массивы, функции;
    \item пользовательские: \texttt{class}, \texttt{struct}, \texttt{enum}, \texttt{union}.
\end{itemize}

По cppreference часто используют более грубое деление:
\begin{itemize}
    \item fundamental;
    \item compound.
\end{itemize}


\subsection{Целочисленные типы данных: размеры, знаковость, диапазоны значений. Псевдонимы фиксированного размера, кодирование, Little Endian.}

К целочисленным типам относятся:
\begin{itemize}
    \item \texttt{char};
    \item \texttt{short};
    \item \texttt{int};
    \item \texttt{long};
    \item \texttt{long long};
\end{itemize}
и их unsigned-варианты.

Размеры зависят от платформы, но для курса важно помнить:
\begin{itemize}
    \item \texttt{char} часто 1 байт;
    \item \texttt{int} часто 4 байта.
\end{itemize}

Псевдонимы фиксированного размера из \texttt{<cstdint>}:
\begin{itemize}
    \item \texttt{std::int8\_t}, \texttt{std::uint8\_t};
    \item \texttt{std::int16\_t}, \texttt{std::uint16\_t};
    \item \texttt{std::int32\_t}, \texttt{std::uint32\_t};
    \item \texttt{std::int64\_t}, \texttt{std::uint64\_t}.
\end{itemize}

Little Endian означает, что младший байт многобайтового числа расположен по меньшему адресу.

\begin{example}
Если \texttt{char c = 'A';} и \texttt{int i = 'A';}, а код символа \texttt{'A'} равен \texttt{65\_10 = 41\_16}, то в памяти:
\begin{itemize}
    \item \texttt{c}: \texttt{41};
    \item \texttt{i}: \texttt{41 00 00 00}.
\end{itemize}
\end{example}


\subsection{Диапазоны значений ц/ч. типов: \texttt{<climits>} и \texttt{<limits>}.}

\texttt{<climits>} даёт макросы:
\begin{itemize}
    \item \texttt{INT\_MIN};
    \item \texttt{INT\_MAX};
    \item и аналогичные для других целых типов.
\end{itemize}

\texttt{<limits>} даёт шаблон:
\begin{verbatim}
std::numeric_limits<T>
\end{verbatim}
через который можно получить свойства типа более общим и ``плюсовым'' способом.


\subsection{Типы чисел с плавающей точкой float и double. Размер и структура.}

Основные типы:
\begin{itemize}
    \item \texttt{float};
    \item \texttt{double};
    \item \texttt{long double}.
\end{itemize}

Обычно:
\begin{itemize}
    \item \texttt{float} занимает 4 байта;
    \item \texttt{double} --- 8 байт.
\end{itemize}

Внутренняя структура обычно включает:
\begin{itemize}
    \item знак;
    \item порядок;
    \item мантиссу.
\end{itemize}


\subsection{Инициализация числовых объектов. Литералы, префиксы/суффиксы. Narrowing.}

Примеры суффиксов:
\begin{itemize}
    \item \texttt{1u};
    \item \texttt{1l};
    \item \texttt{1ll};
    \item \texttt{1.0f};
    \item \texttt{1.0l}.
\end{itemize}

Narrowing --- сужение типа с возможной потерей информации.

Пример:
\begin{verbatim}
int x = 3.14;   // может пройти
int y{3.14};    // ошибка: narrowing
\end{verbatim}

Именно поэтому инициализация через \texttt{\{\}} считается более безопасной.


\subsection{POD-типы, особенность.}

POD --- историческое обозначение ``простых'' типов данных без сложной пользовательской логики.

Для учебного уровня важно помнить:
\begin{itemize}
    \item такие типы имеют простое и предсказуемое представление в памяти;
    \item для них проще рассуждать о побайтном копировании и footprint.
\end{itemize}


\subsection{Выражения и операторы. Классификация операторов, арность, приоритетность. Инструкция-выражение. Различие между инициализацией и присваиванием.}

Выражение --- конструкция, которая вычисляется.

Оператор --- обозначение операции. По арности операторы бывают:
\begin{itemize}
    \item унарные;
    \item бинарные;
    \item тернарные.
\end{itemize}

Приоритет задаёт порядок вычисления.

Инструкция-выражение --- выражение, завершённое \texttt{;}:
\begin{verbatim}
x = y + 1;
f();
\end{verbatim}

Важно различать:
\begin{itemize}
    \item \textbf{инициализация} --- в момент создания объекта;
    \item \textbf{присваивание} --- после того, как объект уже создан.
\end{itemize}


\subsection{Логические выражения и инструкции управления потоком исполнения программы (CF).}

Логические операции:
\begin{itemize}
    \item \texttt{\&\&};
    \item \texttt{||};
    \item \texttt{!};
\end{itemize}

Операторы сравнения:
\begin{itemize}
    \item \texttt{==}, \texttt{!=};
    \item \texttt{<}, \texttt{<=}, \texttt{>}, \texttt{>=}.
\end{itemize}

Основные CF-конструкции:
\begin{itemize}
    \item \texttt{if}, \texttt{if..else};
    \item \texttt{if} с инициализацией;
    \item \texttt{switch}, \texttt{switch} с инициализацией;
    \item тернарный оператор;
    \item \texttt{while}, \texttt{for}, \texttt{range-based for}, \texttt{do..while};
    \item \texttt{break}, \texttt{continue}, \texttt{return}.
\end{itemize}


\subsection{Может ли функция иметь более одной точки выхода с помощью инструкции return? Может ли процедура (void-функция) иметь в теле инструкцию return?}

Да, функция может иметь несколько инструкций \texttt{return}. Это допустимо, хотя в разных стилях программирования к этому относятся по-разному.

Да, \texttt{void}-функция тоже может содержать \texttt{return;}. В таком случае это просто досрочное завершение функции без возвращаемого значения.



\section{Потоки и строки STL}

\subsection{Что подразумевают под моделью ввода-вывода консольного приложения (CLI-application)? Какие основные компоненты ввода-вывода задействуются для него?}

CLI-приложение взаимодействует с пользователем через текстовый ввод и вывод в консоли.

Главные компоненты:
\begin{itemize}
    \item стандартный ввод;
    \item стандартный вывод;
    \item при необходимости --- файловый ввод и вывод.
\end{itemize}

Полезная модель из лекции: консольное приложение можно рассматривать как ``чёрный ящик'', который получает данные через поток ввода и отправляет данные в поток вывода.


\subsection{Что такое стандартный поток ввода и стандартный поток вывода? Какими объектами они представлены? Какая область видимости этих объектов?}

Стандартный поток ввода --- \texttt{std::cin}.  
Стандартный поток вывода --- \texttt{std::cout}.

Это глобальные объекты стандартной библиотеки из пространства имён \texttt{std}.


\subsection{Какой заголовочный файл необходимо подключить, чтобы получить доступ к стандартным потокам ввода/вывода? Где они описаны и где определены?}

Для работы со стандартными потоками подключают:
\begin{verbatim}
#include <iostream>
\end{verbatim}

В заголовке они \textbf{объявлены}, а реальные объекты \textbf{определены} в стандартной библиотеке.


\subsection{Какие типы данных у стандартных потоков ввода и вывода? Какая иерархия наследования для потоковых типов STL?}

\begin{itemize}
    \item \texttt{std::cin} имеет тип \texttt{std::istream};
    \item \texttt{std::cout} имеет тип \texttt{std::ostream}.
\end{itemize}

Более общая иерархия включает:
\begin{itemize}
    \item \texttt{ios\_base};
    \item \texttt{basic\_ios};
    \item \texttt{basic\_istream};
    \item \texttt{basic\_ostream};
    \item \texttt{basic\_iostream};
    \item файловые и строковые производные потоки.
\end{itemize}


\subsection{Какие операции используются для ввода и вывода? Можно ли использовать оператор ввода для потока вывода и наоборот? Как они работают с пробелами? Что возвращают?}

Для ввода используется \texttt{>>}, для вывода --- \texttt{<<}.

\begin{itemize}
    \item \texttt{std::cin >> x;} --- вводит данные из потока;
    \item \texttt{std::cout << x;} --- выводит данные в поток.
\end{itemize}

Оператор \texttt{>>} обычно читает токены, разделённые пробелами.  
Оператор \texttt{<<} выводит данные как есть.

Оба оператора возвращают \textbf{ссылку на поток}, поэтому можно строить цепочки:
\begin{verbatim}
std::cout << a << " " << b;
std::cin >> x >> y;
\end{verbatim}


\subsection{Какие функции ввода и вывода можно использовать наравне с операторами? Как прочитать целиком строку, содержащую пробелы?}

Для чтения целой строки с пробелами используют:
\begin{verbatim}
std::getline(std::cin, s);
\end{verbatim}

Это свободная функция, а не метод строки.

В отличие от \texttt{>>}, она читает строку целиком до символа конца строки.


\subsection{Что такое поток? Почему поток называют абстракцией?}

Поток --- абстракция для ввода или вывода.

Его называют абстракцией, потому что программист работает не с конкретной клавиатурой, экраном или файлом напрямую, а с единым интерфейсом потока символов.


\subsection{Что такое C-строка? Чем она отличается от обычного массива символов? Что такое null terminator?}

C-строка --- это массив символов, который оканчивается \texttt{'\textbackslash 0'}.

Обычный массив символов не обязан содержать завершающий нулевой символ.

\begin{example}
\begin{verbatim}
char a[8] = {'j','u','s','t',' ','a','r','r'};       // не строка
char s[8] = {'c','s','t','r','i','n','g','\0'};      // C-строка
\end{verbatim}
\end{example}

\texttt{null terminator} --- это \textbf{число 0}, а не символ цифры \texttt{'0'}.


\subsection{Что происходит со строковым литералом в C/C++?}

Компилятор превращает строковый литерал в массив символов, то есть в C-строку, размер которой равен числу видимых символов плюс один символ \texttt{'\textbackslash 0'}.

\begin{example}
\begin{verbatim}
char str[] = "Otter";
\end{verbatim}
\end{example}

Здесь реально создаётся массив из шести символов: \texttt{O t t e r \textbackslash 0}.


\subsection{Что такое \texttt{std::string}? Как он связан с \texttt{std::basic\_string<char>}?}

\texttt{std::string} --- стандартный тип C++ для представления строк.

Он является псевдонимом:
\begin{verbatim}
std::string == std::basic_string<char>
\end{verbatim}

Обычно \texttt{std::string} хранит непрерывную последовательность 8-битных символов.


\subsection{Какие базовые операции есть у \texttt{std::string}?}

Часто используемые операции:
\begin{itemize}
    \item \texttt{s.size()}, \texttt{s.length()};
    \item \texttt{s.clear()};
    \item \texttt{s.empty()};
    \item \texttt{s[i]};
    \item \texttt{s + t};
    \item \texttt{s.c\_str()}.
\end{itemize}

\texttt{c\_str()} возвращает указатель на внутренний буфер строки.


\subsection{Как соотносятся размер объекта \texttt{std::string} и длина хранимой в нём строки? Что такое оптимизация коротких строк?}

Размер объекта \texttt{std::string} и длина строки --- это не одно и то же.

Сам объект строки содержит служебные поля, например:
\begin{itemize}
    \item текущую длину;
    \item capacity;
    \item указатель на буфер, либо внутренний локальный буфер.
\end{itemize}

Поэтому \texttt{sizeof(std::string)} не равен длине текста внутри строки.

Оптимизация коротких строк (\textit{SSO, short string optimization}) означает, что короткие строки могут храниться прямо внутри объекта \texttt{std::string}, без отдельной аллокации в куче.


\subsection{Что такое \texttt{std::stringstream} и зачем он нужен?}

\texttt{std::stringstream} --- специальный поток, который сочетает свойства строки и потока.

Он полезен, когда нужно:
\begin{itemize}
    \item разбирать строку как поток;
    \item собирать строку с помощью потоковых операций;
    \item использовать знакомый интерфейс \texttt{>>} и \texttt{<<} для строк.
\end{itemize}



\section{Функциональная декомпозиция, ссылки и константность}

\subsection{Как решить большую задачу так, чтобы полученное решение было корректным, проверяемым, воспроизводимым, переиспользуемым, читаемым?}

Ключевая идея --- \textbf{декомпозиция}.

Большую задачу нужно:
\begin{itemize}
    \item разбить на маленькие кусочки;
    \item каждый кусочек --- ещё на более мелкие подзадачи;
    \item повторять это, пока отдельные части не станут достаточно простыми.
\end{itemize}

Хороший отдельный кусок должен быть:
\begin{itemize}
    \item параметризуемым;
    \item переиспользуемым;
    \item проверяемым отдельно от всей большой программы.
\end{itemize}


\subsection{Что такое функция с точки зрения языка C++? Как определить функцию/процедуру? Какие основные компоненты функции C++?}

Функция --- именованный блок инструкций с интерфейсом и телом.

Основные компоненты:
\begin{itemize}
    \item возвращаемый тип;
    \item имя;
    \item список параметров;
    \item тело функции.
\end{itemize}

Процедура --- это функция с возвращаемым типом \texttt{void}.


\subsection{Что такое неполный тип void? Можно ли породить объект такого типа? а типа void* --- и почему?}

\texttt{void} --- специальный неполный тип, означающий отсутствие значения.

Объект типа \texttt{void} создать нельзя.

А вот \texttt{void*} создать можно, потому что это уже указатель, и он хранит адрес, а не сам объект типа \texttt{void}.


\subsection{Что такое свободная функция (non-member function) и что такое функция-член (member function)? Привести примеры.}

Свободная функция не принадлежит классу:
\begin{verbatim}
int sum(int a, int b)
{
    return a + b;
}
\end{verbatim}

Функция-член принадлежит классу:
\begin{verbatim}
struct Point
{
    double x, y;
    double norm() const { return x * x + y * y; }
};
\end{verbatim}


\subsection{Является ли функция \texttt{std::getline} свободной или функцией-членом?}

\texttt{std::getline} --- \textbf{свободная функция}.


\subsection{Является ли статическая функция (помечена ключевым словом static) \texttt{bar()}, объявленная в классе Foo, свободной или функцией-членом? Как вызвать такую функцию?}

Если функция объявлена внутри класса как \texttt{static}, то это всё равно \textbf{функция-член}, но без неявного параметра \texttt{this}.

Вызывается так:
\begin{verbatim}
Foo::bar();
\end{verbatim}


\subsection{Что такое side-effect при вызове свободных функций и почему его наличие считается нежелательным? Как с помощью static можно внести в функцию некоторое состояние?}

Side-effect --- это изменение состояния, выходящее за рамки ``чистого'' вычисления результата функции.

Он часто считается нежелательным, потому что:
\begin{itemize}
    \item ухудшает предсказуемость;
    \item усложняет тестирование;
    \item ухудшает переиспользуемость.
\end{itemize}

Пример внутреннего состояния функции:
\begin{verbatim}
int f()
{
    static int counter = 0;
    return ++counter;
}
\end{verbatim}

Здесь \texttt{static} сохраняет значение между вызовами и тем самым создаёт side-effect.


\subsection{Что такое методы для классов и структур? Приведите пример метода для класса \texttt{std::string}.}

Методы --- это функции-члены классов и структур.

Примеры методов \texttt{std::string}:
\begin{itemize}
    \item \texttt{size()};
    \item \texttt{empty()};
    \item \texttt{clear()};
    \item \texttt{c\_str()}.
\end{itemize}


\subsection{Что такое список формальных параметров функции? Можно ли составить список ф.п., состоящий только из типов?}

Список формальных параметров --- это перечень параметров в заголовке функции.

В определении функции нужны имена параметров, если они используются в теле.

В прототипе имена можно опускать:
\begin{verbatim}
int sum(int, int);
\end{verbatim}

То есть список может состоять только из типов, если речь идёт о прототипе.


\subsection{Что такое сигнатура функции? Для чего необходимо рассматривать сигнатуру отдельно от заголовка функции?}

Сигнатура функции --- это часть её интерфейса, которая отличает одну перегрузку от другой.

Обычно в неё входят:
\begin{itemize}
    \item имя функции;
    \item типы параметров;
    \item в C++ также qualifiers, если они участвуют в различении перегрузок.
\end{itemize}

Сигнатуру рассматривают отдельно, потому что именно она важна для перегрузки, а не весь текст заголовка целиком.


\subsection{Чем различаются заголовок функции от прототипа функции? Что из них относится к интерфейсной части функции?}

Заголовок функции --- это строка с возвращаемым типом, именем и параметрами.

Если к заголовку добавить \texttt{;} и не писать тело, получится \textbf{прототип функции}.

Именно прототип относится к интерфейсной части функции в явном виде.


\subsection{Параметры функции: передача по значению и по ссылке. Когда, зачем и почему?}

Передача по значению означает копирование аргумента в параметр функции.

Передача по ссылке означает работу с исходным объектом без копии.

Идиоматические правила:
\begin{itemize}
    \item большие объекты для чтения --- через \texttt{const T\&};
    \item маленькие объекты для чтения --- по значению;
    \item объекты для чтения и записи --- через \texttt{T\&}.
\end{itemize}


\subsection{Что такое константная ссылка?}

Константная ссылка --- это ещё одно имя для объекта, через которое можно только читать, но нельзя изменять объект.

\begin{example}
\begin{verbatim}
const std::string& s
\end{verbatim}
\end{example}

Полезная интуиция из лекции: это однонаправленный ``портал'' только для чтения.


\subsection{Идиоматический приём: константная ссылка в range-based for.}

Если элементы диапазона большие, лучше итерироваться так:
\begin{verbatim}
for (const auto& el : container)
{
    ...
}
\end{verbatim}

Это позволяет:
\begin{itemize}
    \item не копировать каждый элемент;
    \item не давать возможности случайно изменить элемент;
    \item сохранить эффективность.
\end{itemize}


\subsection{«Большие» и «малые» объекты: правило трёх указателей.}

Практический критерий:
\begin{itemize}
    \item если объект примерно не больше трёх указателей, его считают ``небольшим'';
    \item если существенно больше --- ``большим''.
\end{itemize}

Это не закон стандарта, а инженерное идиоматическое правило курса для выбора способа передачи объектов в функции.


\subsection{С какой целью функция, получающая параметр по ссылке, может возвращать тот же параметр тоже по ссылке (pass-through)?}

Это приём pass-through.

Он нужен, чтобы:
\begin{itemize}
    \item не создавать копию;
    \item сохранить ссылочную семантику;
    \item поддержать цепочки вызовов.
\end{itemize}

\begin{example}
\begin{verbatim}
int& pass(int& x)
{
    return x;
}
\end{verbatim}
\end{example}



\section{Стандартная библиотека: основные концепции и линейные контейнеры}

\subsection{Что такое стандартная библиотека C++ и зачем она нужна?}

Стандартная библиотека --- часть стандарта C++, предоставляющая готовые типы, алгоритмы, потоки, контейнеры и другие средства общего назначения.

Она нужна, чтобы не реализовывать заново базовые вещи: строки, потоки, контейнеры, итераторы, алгоритмы и т.д.


\subsection{Что такое контейнер STL? Какие большие группы контейнеров выделяют?}

Контейнер STL --- тип данных, предназначенный для хранения коллекции элементов.

В лекциях выделялись три большие группы:
\begin{itemize}
    \item линейные контейнеры;
    \item ассоциативные контейнеры;
    \item неупорядоченные контейнеры.
\end{itemize}


\subsection{Какие контейнеры относятся к линейным?}

К линейным контейнерам относились:
\begin{itemize}
    \item \texttt{std::array};
    \item \texttt{std::vector};
    \item \texttt{std::deque};
    \item \texttt{std::forward\_list};
    \item \texttt{std::list}.
\end{itemize}


\subsection{Что такое \texttt{std::vector}? В чём его главные свойства?}

\texttt{std::vector} --- линейный контейнер динамического размера.

Главные свойства:
\begin{itemize}
    \item хранит элементы подряд в памяти;
    \item даёт быстрый доступ по индексу;
    \item поддерживает рост контейнера;
    \item при нехватке памяти может перевыделять внутренний буфер.
\end{itemize}


\subsection{Чем \texttt{std::vector} отличается от C-массива и от \texttt{std::array}?}

Общее с массивами:
\begin{itemize}
    \item последовательное расположение в памяти;
    \item быстрый доступ по индексу.
\end{itemize}

Отличия:
\begin{itemize}
    \item C-массив не знает своего размера как объектного интерфейса;
    \item \texttt{std::array} имеет фиксированный размер;
    \item \texttt{std::vector} может динамически расти.
\end{itemize}


\subsection{Что такое size и capacity у \texttt{std::vector}?}

\begin{itemize}
    \item \texttt{size()} --- текущее число элементов;
    \item \texttt{capacity()} --- объём памяти, уже зарезервированный под элементы без нового перевыделения.
\end{itemize}

При добавлении элементов capacity растёт не по одному, а обычно скачками.


\subsection{Для чего нужен \texttt{reserve()}?}

\texttt{reserve(n)} заранее резервирует память под не менее чем \texttt{n} элементов.

Это полезно, чтобы:
\begin{itemize}
    \item уменьшить число перевыделений памяти;
    \item не инвалидировать итераторы лишний раз;
    \item ускорить массовое добавление элементов.
\end{itemize}


\subsection{Что такое матрица в контексте STL-решений курса?}

В курсе матрицу удобно представлять как ``вектор векторов'':
\begin{verbatim}
std::vector<std::vector<int>>
\end{verbatim}

Это естественное и простое представление двумерной структуры через линейные контейнеры.


\subsection{Что такое алиас типа? Чем \texttt{using} отличается от \texttt{typedef}?}

Алиас типа --- новое имя для уже существующего типа.

Через \texttt{typedef}:
\begin{verbatim}
typedef std::vector<int> Row;
\end{verbatim}

Через \texttt{using}:
\begin{verbatim}
using Row = std::vector<int>;
\end{verbatim}

Обе записи делают одно и то же, но \texttt{using} обычно считается более современным и удобным.


\subsection{Приведите пример алиасов для ряда и зрительного зала.}

\begin{verbatim}
using Row = std::vector<int>;
using SittingPlan = std::vector<Row>;
\end{verbatim}



\section{Стандартная библиотека: итераторы}

\subsection{Что такое итератор?}

Итератор --- объект, который может итерировать элементы некоторой коллекции, то есть перемещаться от элемента к элементу.

Он задаёт определённую позицию в контейнере и служит мостом между контейнерами и алгоритмами.


\subsection{Что означают \texttt{begin()} и \texttt{end()}? Почему интервал записывают как \texttt{[begin, end)}?}

\begin{itemize}
    \item \texttt{begin()} --- итератор на первый элемент;
    \item \texttt{end()} --- итератор на позицию \emph{за последним} элементом.
\end{itemize}

Интервал записывают как \texttt{[begin, end)}, то есть левая граница включена, правая не включена.

Поэтому:
\begin{itemize}
    \item \texttt{*begin()} обычно допустимо для непустого контейнера;
    \item \texttt{*end()} --- недопустимо.
\end{itemize}


\subsection{Можно ли сделать копию итератора \texttt{c.begin()}? Можно ли эту копию разыменовать? Можно ли эту копию двигать вперёд?}

Да:
\begin{itemize}
    \item копию сделать можно;
    \item разыменовать обычно можно, если контейнер не пуст;
    \item двигать вперёд можно, если это допускает категория итератора.
\end{itemize}


\subsection{Можно ли сделать копию итератора \texttt{c.end()}? Можно ли эту копию разыменовать? Можно ли эту копию двигать назад?}

Копию сделать можно.

Разыменовывать \texttt{end()} нельзя.

Двигать назад можно не всегда, а только если категория итератора не ниже bidirectional и если это осмысленно для данного контейнера и позиции.


\subsection{Какой вариант сравнения итераторов при переборе коллекции следует предпочесть: \texttt{it != end()} или \texttt{it < end()}? Почему?}

Следует предпочитать:
\begin{verbatim}
it != end()
\end{verbatim}

Причина: сравнение \texttt{<} поддерживают не все итераторы. Оно доступно только для random access итераторов. А \texttt{!=} подходит гораздо более широко.


\subsection{Какие категории итераторов выделяли до стандарта C++20? Чем они различаются? Как категории соотносятся друг с другом?}

В упрощённом виде в лекции разбирались:
\begin{itemize}
    \item \textbf{forward};
    \item \textbf{bidirectional};
    \item \textbf{random access}.
\end{itemize}

Соотношение такое:
\begin{itemize}
    \item random access включает возможности bidirectional;
    \item bidirectional включает возможности forward.
\end{itemize}

Различаются набором допустимых операций:
\begin{itemize}
    \item forward: \texttt{++it}, \texttt{it++};
    \item bidirectional: ещё и \texttt{--it}, \texttt{it--};
    \item random access: ещё и \texttt{it + n}, \texttt{it - n}, сравнение порядка и т.д.
\end{itemize}


\subsection{При разработке нового алгоритма какую категорию итераторов следует требовать?}

Правильная политика: требовать \textbf{минимальную} категорию итераторов, которая действительно нужна алгоритму.

Это делает алгоритм применимым к большему числу контейнеров.


\subsection{Что происходит с итератором после удаления элемента из коллекции? Как следует поступить с этим итератором после удаления элемента?}

После удаления элемента итератор на этот элемент обычно становится \textbf{недействительным}.

Общая безопасная политика:
\begin{itemize}
    \item не использовать старый инвалидированный итератор;
    \item брать новый итератор из возвращаемого значения операции удаления, если контейнер это поддерживает.
\end{itemize}



\section{Косвенная адресация}

\subsection{Что такое косвенная адресация? Какие основные варианты косвенности разбирались в курсе?}

Косвенная адресация --- это обращение к объекту не напрямую, а через некоторую ``прослойку''.

В курсе разбирались:
\begin{itemize}
    \item указатели;
    \item ссылки;
    \item итераторы;
    \item \texttt{std::string\_view}.
\end{itemize}


\subsection{В чём общая идея указателя, ссылки и итератора как средств косвенности?}

Во всех трёх случаях мы работаем не с объектом ``сам по себе'', а через сущность, которая позволяет к нему обратиться.

Различия:
\begin{itemize}
    \item указатель требует явного разыменования;
    \item ссылка ведёт себя почти как само имя объекта;
    \item итератор ведёт к элементу контейнера и умеет ``двигаться'' по коллекции.
\end{itemize}


\subsection{Чем указатель отличается от ссылки?}

Указатель:
\begin{itemize}
    \item является отдельным объектом;
    \item хранит адрес;
    \item может быть перенаправлен;
    \item требует разыменования \texttt{*} или \texttt{->}.
\end{itemize}

Ссылка:
\begin{itemize}
    \item обычно рассматривается как ещё одно имя объекта;
    \item после инициализации не перепривязывается;
    \item не требует явного разыменования.
\end{itemize}


\subsection{Что значит левоконстантный, правоконстантный и дваждыконстантный указатель?}

\begin{itemize}
    \item \texttt{const int* p;} --- константен объект по адресу, но не сам указатель;
    \item \texttt{int* const p = ...;} --- константен сам указатель, но объект по адресу можно менять;
    \item \texttt{const int* const p = ...;} --- константны и указатель, и объект по адресу.
\end{itemize}

Именно это в курсе и называли различными разновидностями константности указателя.


\subsection{Что можно сказать о константной и неконстантной ссылке?}

\begin{itemize}
    \item \texttt{T\&} --- неконстантная ссылка, через неё можно менять объект;
    \item \texttt{const T\&} --- константная ссылка, через неё можно только читать объект.
\end{itemize}


\subsection{Что такое \texttt{const\_iterator}, \texttt{const iterator} и \texttt{const const\_iterator}?}

\begin{itemize}
    \item \texttt{const\_iterator} --- итератор, который не позволяет изменять элемент коллекции;
    \item \texttt{const iterator} --- константный объект-итератор, который нельзя ``передвигать'', но тип элемента может зависеть от самого итератора;
    \item \texttt{const const\_iterator} --- одновременно константен сам объект-итератор и запрещено менять элемент через него.
\end{itemize}

Здесь важно различать:
\begin{itemize}
    \item константность \textbf{объекта};
    \item константность \textbf{косвенности}.
\end{itemize}



\section{Класс std::string\_view}

\subsection{Что такое \texttt{std::string\_view}?}

\texttt{std::string\_view} --- это легковесная невладеющая обёртка над строкой или подстрокой.

Он не владеет символами, а только ``смотрит'' на уже существующий диапазон символов.


\subsection{Почему \texttt{std::string\_view} называют невладеющим типом?}

Потому что он не управляет временем жизни символов, на которые указывает.

Следовательно, строка, на которую смотрит \texttt{string\_view}, должна жить дольше самого \texttt{string\_view}.


\subsection{Почему \texttt{std::string\_view} удобно передавать по значению?}

Потому что это лёгкий тип: обычно он содержит только указатель и длину.

Из-за этого его копирование дешёвое, и правило передачи ``малых'' объектов по значению здесь подходит естественно.


\subsection{Можно ли модифицировать строку через \texttt{std::string\_view}?}

Нет. В лекциях подчёркивалось, что \texttt{std::string\_view} --- \textbf{immutable} и допускает только немодифицирующие операции над строкой.


\subsection{Гарантирует ли \texttt{std::string\_view} наличие \texttt{'\textbackslash 0'} в конце?}

Нет, не гарантирует.

Поэтому его неудобно напрямую использовать с legacy C-кодом, который ожидает классическую C-строку.


\subsection{В каком случае \texttt{std::string\_view} особенно хорош?}

Когда нужно эффективно исследовать строки и подстроки без копирования и без модификации исходных данных.


\subsection{Как \texttt{std::string\_view} связан с идеей косвенной адресации?}

Он является ещё одним примером косвенности: объект \texttt{string\_view} сам текст не хранит, а лишь косвенно обращается к уже существующим символам.



\section{Ассоциативные контейнеры std::set и std::map}

\subsection{Что такое \texttt{std::set}?}

\texttt{std::set} --- ассоциативный контейнер, который хранит элементы без дубликатов.

Элементы автоматически упорядочиваются согласно критерию сравнения.


\subsection{Что такое \texttt{std::map}?}

\texttt{std::map} --- ассоциативный контейнер, хранящий пары ``ключ--значение''.

Его можно рассматривать как ассоциативный массив.


\subsection{Что является элементом \texttt{std::map}?}

Элемент \texttt{std::map} представляется объектом типа \texttt{std::pair<Key, Value>}.


\subsection{Какой тип обычно используют для представления пары объектов?}

Для этого используют служебный тип:
\begin{verbatim}
std::pair<Type1, Type2>
\end{verbatim}

У него есть поля:
\begin{itemize}
    \item \texttt{first};
    \item \texttt{second}.
\end{itemize}


\subsection{Какая типичная внутренняя реализация у \texttt{std::set} и \texttt{std::map}?}

Типичная реализация --- сбалансированное бинарное дерево поиска, обычно красно-чёрное дерево.


\subsection{Какова асимптотика основных операций у \texttt{std::set} и \texttt{std::map}?}

Для основных операций:
\begin{itemize}
    \item поиск --- \texttt{O(log n)};
    \item вставка --- \texttt{O(log n)};
    \item удаление --- \texttt{O(log n)}.
\end{itemize}

Полный обход всех элементов --- \texttt{O(n)}.


\subsection{Почему при поиске и вставке в \texttt{set}/\texttt{map} используется сравнение?}

Потому что контейнеры организованы как дерево поиска, и на каждом шаге нужно решать, идти ``влево'' или ``вправо''.

Для этого и нужен критерий ``меньше''.


\subsection{Почему для \texttt{std::set} обычно достаточно \texttt{operator<}?}

Потому что стандартный компаратор по умолчанию --- \texttt{std::less<Key>}, а он, в простейшем случае, использует именно \texttt{operator<}.


\subsection{Можно ли хранить в \texttt{set} дубликаты?}

Нет, \texttt{std::set} хранит элементы без дубликатов.

Если нужны повторяющиеся ключи, используют \texttt{std::multiset}.


\subsection{Можно ли хранить в \texttt{map} несколько одинаковых ключей?}

Нет, в \texttt{std::map} ключи уникальны.

Если нужны повторяющиеся ключи, используют \texttt{std::multimap}.



\section{Ассоциативные контейнеры, компараторы, сравнение и сортировка}

\subsection{Каково назначение параметра \texttt{Compare} в шаблонах \texttt{std::set<Key, Compare>} и \texttt{std::map<Key, Value, Compare>}?}

Параметр \texttt{Compare} задаёт правило упорядочивания элементов внутри контейнера.

Именно через него контейнер понимает, какой из двух ключей ``меньше'' и как размещать элементы в дереве.


\subsection{Что такое компаратор?}

Компаратор --- объект или функция, задающая отношение порядка между двумя значениями.

В простейшем случае он отвечает на вопрос: должно ли первое значение идти раньше второго.


\subsection{Что сравнивает компаратор у \texttt{std::set} и у \texttt{std::map}?}

\begin{itemize}
    \item у \texttt{std::set<Key, Compare>} компаратор сравнивает элементы типа \texttt{Key};
    \item у \texttt{std::map<Key, Value, Compare>} компаратор сравнивает только \textbf{ключи} типа \texttt{Key}.
\end{itemize}


\subsection{Что такое \texttt{std::less}?}

\texttt{std::less<T>} --- стандартный компаратор по умолчанию.

В простейшем варианте он использует \texttt{operator<}:
\begin{verbatim}
lhs < rhs
\end{verbatim}


\subsection{Можно ли задать свой компаратор?}

Да.

Это можно сделать:
\begin{itemize}
    \item через функциональный объект;
    \item через тип указателя на функцию;
    \item через лямбда-выражение, если это поддержано выбранной формой записи.
\end{itemize}


\subsection{Что такое свободная функция в роли компаратора?}

Это обычная функция, которую можно использовать для сравнения двух значений.

Например, тип указателя на такую функцию можно описать так:
\begin{verbatim}
using MyCompareFunPtr = bool(*)(int, int);
\end{verbatim}


\subsection{Почему компаратор должен задавать устойчивое и разумное правило порядка?}

Потому что от него зависит структура дерева, а значит:
\begin{itemize}
    \item корректность поиска;
    \item корректность вставки;
    \item отсутствие логических противоречий в порядке элементов.
\end{itemize}

Если правило сравнения ``ломаное'', контейнер ведёт себя непредсказуемо.


\subsection{Какая типичная внутренняя реализация у \texttt{std::map} и какова асимптотика операций?}

Как и \texttt{std::set}, \texttt{std::map} обычно реализуется как сбалансированное бинарное дерево поиска, чаще всего красно-чёрное дерево.

Отсюда:
\begin{itemize}
    \item поиск --- \texttt{O(log n)};
    \item вставка --- \texttt{O(log n)};
    \item удаление --- \texttt{O(log n)};
    \item обход --- \texttt{O(n)}.
\end{itemize}


\subsection{Можно ли использовать один и тот же компаратор в разных контейнерах?}

Да, если этот компаратор умеет сравнивать те значения, которые надо сравнивать в данном контейнере.

Например:
\begin{itemize}
    \item для \texttt{std::set<T>} он сравнивает элементы типа \texttt{T};
    \item для \texttt{std::map<T, U>} --- ключи типа \texttt{T};
    \item для \texttt{std::vector<T>} при сортировке --- элементы типа \texttt{T}.
\end{itemize}

Главное, чтобы смысл порядка был согласован с задачей.
